\documentclass{amsart}
% For dealing with references we use the comment environment
\usepackage{verbatim}
\newenvironment{reference}{\comment}{\endcomment}
%\newenvironment{reference}{}{}
\newenvironment{slogan}{\comment}{\endcomment}
\newenvironment{history}{\comment}{\endcomment}

% For commutative diagrams we use Xy-pic
\usepackage[all]{xy}

% We use 2cell for 2-commutative diagrams.
\xyoption{2cell}
\UseAllTwocells

% We use multicol for the list of chapters between chapters
\usepackage{multicol}

% This is generally recommended for better output
\usepackage{lmodern}
\usepackage[T1]{fontenc}

% For cross-file-references

% Package for hypertext links:
\usepackage{hyperref}

% For any local file, say "hello.tex" you want to link to please

% Theorem environments.
%
\theoremstyle{plain}
\newtheorem{theorem}[subsection]{Theorem}
\newtheorem{proposition}[subsection]{Proposition}
\newtheorem{lemma}[subsection]{Lemma}

\theoremstyle{definition}
\newtheorem{definition}[subsection]{Definition}
\newtheorem{example}[subsection]{Example}
\newtheorem{exercise}[subsection]{Exercise}
\newtheorem{situation}[subsection]{Situation}

\theoremstyle{remark}
\newtheorem{remark}[subsection]{Remark}
\newtheorem{remarks}[subsection]{Remarks}

\numberwithin{equation}{subsection}

% Macros
%
\def\lim{\mathop{\mathrm{lim}}\nolimits}
\def\colim{\mathop{\mathrm{colim}}\nolimits}
\def\Spec{\mathop{\mathrm{Spec}}}
\def\Hom{\mathop{\mathrm{Hom}}\nolimits}
\def\Ext{\mathop{\mathrm{Ext}}\nolimits}
\def\SheafHom{\mathop{\mathcal{H}\!\mathit{om}}\nolimits}
\def\SheafExt{\mathop{\mathcal{E}\!\mathit{xt}}\nolimits}
\def\Sch{\mathit{Sch}}
\def\Mor{\mathop{\mathrm{Mor}}\nolimits}
\def\Ob{\mathop{\mathrm{Ob}}\nolimits}
\def\Sh{\mathop{\mathit{Sh}}\nolimits}
\def\NL{\mathop{N\!L}\nolimits}
\def\CH{\mathop{\mathrm{CH}}\nolimits}
\def\proetale{{pro\text{-}\acute{e}tale}}
\def\etale{{\acute{e}tale}}
\def\QCoh{\mathit{QCoh}}
\def\Ker{\mathop{\mathrm{Ker}}}
\def\Im{\mathop{\mathrm{Im}}}
\def\Coker{\mathop{\mathrm{Coker}}}
\def\Coim{\mathop{\mathrm{Coim}}}

% Boxtimes
%
\DeclareMathSymbol{\boxtimes}{\mathbin}{AMSa}{"02}

%
% Macros for moduli stacks/spaces
%
\def\QCohstack{\mathcal{QC}\!\mathit{oh}}
\def\Cohstack{\mathcal{C}\!\mathit{oh}}
\def\Spacesstack{\mathcal{S}\!\mathit{paces}}
\def\Quotfunctor{\mathrm{Quot}}
\def\Hilbfunctor{\mathrm{Hilb}}
\def\Curvesstack{\mathcal{C}\!\mathit{urves}}
\def\Polarizedstack{\mathcal{P}\!\mathit{olarized}}
\def\Complexesstack{\mathcal{C}\!\mathit{omplexes}}
% \Pic is the operator that assigns to X its picard group, usage \Pic(X)
% \Picardstack_{X/B} denotes the Picard stack of X over B
% \Picardfunctor_{X/B} denotes the Picard functor of X over B
\def\Pic{\mathop{\mathrm{Pic}}\nolimits}
\def\Picardstack{\mathcal{P}\!\mathit{ic}}
\def\Picardfunctor{\mathrm{Pic}}
\def\Deformationcategory{\mathcal{D}\!\mathit{ef}}

\setlength{\emergencystretch}{3em}
\begin{document}
\title[Stacks Project, Tag 0AIC]{722 --- Affine formal spaces arise from weakly admissible rings exactly when limit projections surject}
\maketitle
\noindent Copyright (C) 2005--2025 Johan de Jong. Stacks Project authors. Source: \href{https://stacks.math.columbia.edu/tag/0AIC}{Tag 0AIC}.

\noindent Editorial difficulty note (not author text). Difficulty H2: The author recovers the weak ideals of definition from thickenings, constructs a complete limit topology, and proves presentation independence through cofinality. This formal-geometric reconstruction theorem is beyond ordinary qualification material..

\noindent Editorial provenance note (not author text). History: excerpt from revision \texttt{a0444\allowbreak 6e57e\allowbreak c1fbc\allowbreak 252a8\allowbreak 71afc\allowbreak ec775\allowbreak 2fb28\allowbreak 07b14}, formal-spaces.tex, lines 1917--2016. Reference presentation was adapted; original-author context and core support blocks were retained or added. The mathematical wording is unchanged. Licensed under GNU FDL 1.2 or later; no invariant sections or cover texts. See ../licenses/COPYING. Context blocks reproduce author definitions and supporting results; they are not additional counted problems.

\begin{definition}
\label{definition-affine-formal-algebraic-space}
Let $S$ be a scheme. We say a sheaf $X$ on $(\Sch/S)_{fppf}$ is an
{\it affine formal algebraic space} if there exist
\begin{enumerate}
\item a directed set $\Lambda$,
\item a system $(X_\lambda, f_{\lambda \mu})$ over $\Lambda$
in $(\Sch/S)_{fppf}$ where
\begin{enumerate}
\item each $X_\lambda$ is affine,
\item each $f_{\lambda \mu} : X_\lambda \to X_\mu$ is a thickening,
\end{enumerate}
\end{enumerate}
such that
$$
X \cong \colim_{\lambda \in \Lambda} X_\lambda
$$
as fppf sheaves and $X$ satisfies a set theoretic condition
(see Remark \ref{remark-set-theoretic}). A
{\it morphism of affine formal algebraic spaces}
over $S$ is a map of sheaves.
\end{definition}

\begin{remark}[McQuillan's variant]
\label{remark-mcquillan}
There is a variant of the construction of formal schemes due to
McQuillan, see \href{https://stacks.math.columbia.edu/bibliography}{Bibliography: McQuillan}.
He suggests a slight weakening of the condition of admissibility.
Namely, recall that an admissible topological ring is a complete
(and separated by our conventions) topological ring $A$
which is linearly topologized such that there exists an
ideal of definition: an
open ideal $I$ such that any neighbourhood of $0$ contains $I^n$
for some $n \geq 1$.
McQuillan works with what we will call {\it weakly admissible}
topological rings. A weakly admissible topological ring $A$ is a
complete (and separated by our conventions) topological ring
which is linearly topologized such that there exists an
{\it weak ideal of definition}: an open ideal $I$ such that
for all $f \in I$ we have
$f^n \to 0$ for $n \to \infty$. Similarly to the admissible case,
if $I$ is a weak ideal of definition and $J \subset A$ is an
open ideal, then $I \cap J$ is a weak ideal of definition.
Thus the weak ideals of definition form a fundamental system of
open neighbourhoods of $0$ and
one can proceed along much the same route as above
to define a larger category of formal schemes based
on this notion. The analogues of Lemmas \href{https://stacks.math.columbia.edu/tag/0AI1}{lemma fully faithful (Tag 0AI1)} and
\href{https://stacks.math.columbia.edu/tag/0AI2}{lemma formal scheme sheaf fppf (Tag 0AI2)}
still hold in this setting (with the same proof).
\end{remark}

\begin{remark}
\label{remark-set-theoretic}
Let $S$ be a scheme and let $(\Sch/S)_{fppf}$ be a big fppf site as
in Topologies, Definition \ref{topologies-definition-big-small-fppf}.
As our set theoretic condition on $X$ in
Definitions \ref{definition-affine-formal-algebraic-space} and
\ref{definition-formal-algebraic-space} we take:
there exist objects $U, R$ of $(\Sch/S)_{fppf}$, a
morphism $U \to X$ which is a surjection of fppf sheaves, and
a morphism $R \to U \times_X U$ which is a surjection of fppf sheaves.
In other words, we require our sheaf to be a coequalizer of
two maps between representable sheaves.
Here are some observations which imply this notion behaves
reasonably well:
\begin{enumerate}
\item Suppose $X = \colim_{\lambda \in \Lambda} X_\lambda$
and the system satisfies conditions (1) and (2) of
Definition \ref{definition-affine-formal-algebraic-space}. Then
$U = \coprod_{\lambda \in \Lambda} X_\lambda \to X$ is a surjection
of fppf sheaves. Moreover, $U \times_X U$ is a closed subscheme
of $U \times_S U$ by Lemma \href{https://stacks.math.columbia.edu/tag/0AI8}{lemma diagonal affine formal algebraic space (Tag 0AI8)}.
Hence if $U$ is representable by an object of $(\Sch/S)_{fppf}$
then $U \times_S U$ is too (see Sets, Lemma \href{https://stacks.math.columbia.edu/tag/000R}{lemma what is in it (Tag 000R)})
and the set theoretic condition is satisfied. This is always the case
if $\Lambda$ is countable, see Sets, Lemma \href{https://stacks.math.columbia.edu/tag/000R}{lemma what is in it (Tag 000R)}.
\item Sanity check. Let $\{X_i \to X\}_{i \in I}$ be as in
Definition \ref{definition-formal-algebraic-space}
(with the set theoretic condition as formulated above)
and assume that each $X_i$ is actually an affine scheme.
Then $X$ is an algebraic space. Namely, if we choose a larger
big fppf site $(\Sch'/S)_{fppf}$ such that $U' = \coprod X_i$
and $R' = \coprod X_i \times_X X_j$ are representable by objects
in it, then $X' = U'/R'$ will be an object of the category
of algebraic spaces for this choice. Then an application of
Spaces, Lemma \href{https://stacks.math.columbia.edu/tag/04W1}{lemma fully faithful (Tag 04W1)} shows that
$X$ is an algebraic space for $(\Sch/S)_{fppf}$.
\item Let $\{X_i \to X\}_{i \in I}$ be a family of maps of sheaves
satisfying conditions (1), (2), (3) of
Definition \ref{definition-formal-algebraic-space}.
For each $i$ we can pick $U_i \in \Ob((\Sch/S)_{fppf})$
and $U_i \to X_i$ which is a surjection of sheaves.
Thus if $I$ is not too large (for example countable) then
$U = \coprod U_i \to X$ is a surjection of sheaves and
$U$ is representable by an object of $(\Sch/S)_{fppf}$.
To get $R \in \Ob((\Sch/S)_{fppf})$ surjecting onto $U \times_X U$
it suffices to assume the diagonal $\Delta : X \to X \times_S X$ is not
too wild, for example this always works if the diagonal of $X$ is
quasi-compact, i.e., $X$ is quasi-separated.
\end{enumerate}
\end{remark}

\begin{definition}
\label{topologies-definition-big-small-fppf}
Let $S$ be a scheme. Let $\Sch_{fppf}$ be a big fppf
site containing $S$.
\begin{enumerate}
\item The {\it big fppf site of $S$}, denoted
$(\Sch/S)_{fppf}$, is the site $\Sch_{fppf}/S$
introduced in Sites, Section \href{https://stacks.math.columbia.edu/tag/00XZ}{section localize (Tag 00XZ)}.
\item The {\it big affine fppf site of $S$}, denoted
$(\textit{Aff}/S)_{fppf}$, is the full subcategory of
$(\Sch/S)_{fppf}$ whose objects are affine $U/S$.
A covering of $(\textit{Aff}/S)_{fppf}$ is any covering
$\{U_i \to U\}$ of $(\Sch/S)_{fppf}$ which is a
standard fppf covering.
\end{enumerate}
\end{definition}

\begin{definition}
\label{definition-formal-algebraic-space}
Let $S$ be a scheme. We say a sheaf $X$ on $(\Sch/S)_{fppf}$ is a
{\it formal algebraic space} if there exist a family of maps
$\{X_i \to X\}_{i \in I}$ of sheaves such that
\begin{enumerate}
\item $X_i$ is an affine formal algebraic space,
\item $X_i \to X$ is representable by algebraic spaces and \'etale,
\item $\coprod X_i \to X$ is surjective as a map of sheaves
\end{enumerate}
and $X$ satisfies a set theoretic condition
(see Remark \ref{remark-set-theoretic}). A
{\it morphism of formal algebraic spaces}
over $S$ is a map of sheaves.
\end{definition}

\begin{lemma}
\label{lemma-mcquillan-affine-formal-algebraic-space}
Let $S$ be a scheme. Let $X$ be an fppf sheaf on $(\Sch/S)_{fppf}$
which satisfies the set theoretic condition of
Remark \ref{remark-set-theoretic}.
The following are equivalent:
\begin{enumerate}
\item there exists a weakly admissible topological ring $A$ over $S$
(see Remark \ref{remark-mcquillan}) such that
$X = \colim_{I \subset A\text{ weak ideal of definition}} \Spec(A/I)$,
\item $X$ is an affine formal algebraic space and
there exists an $S$-algebra $A$ and a map $X \to \Spec(A)$
such that for a closed immersion $T \to X$ with $T$ an affine scheme
the composition $T \to \Spec(A)$ is a closed immersion,
\item $X$ is an affine formal algebraic space and
there exists an $S$-algebra $A$ and a map $X \to \Spec(A)$
such that for a closed immersion $T \to X$ with $T$ a scheme
the composition $T \to \Spec(A)$ is a closed immersion,
\item $X$ is an affine formal algebraic space and
for some choice of $X = \colim X_\lambda$ as in
Definition \ref{definition-affine-formal-algebraic-space}
the projections $\lim \Gamma(X_\lambda, \mathcal{O}_{X_\lambda})
\to \Gamma(X_\lambda, \mathcal{O}_{X_\lambda})$ are surjective,
\item $X$ is an affine formal algebraic space and for any choice
of $X = \colim X_\lambda$ as in
Definition \ref{definition-affine-formal-algebraic-space}
the projections $\lim \Gamma(X_\lambda, \mathcal{O}_{X_\lambda})
\to \Gamma(X_\lambda, \mathcal{O}_{X_\lambda})$ are surjective.
\end{enumerate}
Moreover, the weakly admissible topological ring is
$A = \lim \Gamma(X_\lambda, \mathcal{O}_{X_\lambda})$
endowed with its limit topology and the weak ideals of definition
classify exactly the morphisms $T \to X$ which are representable
and thickenings.
\end{lemma}

\begin{proof}
It is clear that (5) implies (4).

\medskip\noindent
Assume (4) for $X = \colim X_\lambda$ as in
Definition \ref{definition-affine-formal-algebraic-space}.
Set $A = \lim \Gamma(X_\lambda, \mathcal{O}_{X_\lambda})$.
Let $T \to X$ be a closed immersion with $T$ a scheme
(note that $T \to X$ is representable by
Lemma \href{https://stacks.math.columbia.edu/tag/0AI8}{lemma diagonal affine formal algebraic space (Tag 0AI8)}).
Since $X_\lambda \to X$ is a thickening, so is
$X_\lambda \times_X T \to T$. On the other hand,
$X_\lambda \times_X T \to X_\lambda$ is a closed immersion,
hence $X_\lambda \times_X T$ is affine. Hence $T$ is affine
by Limits, Proposition \href{https://stacks.math.columbia.edu/tag/05YU}{proposition affine (Tag 05YU)}.
Then $T \to X$ factors through $X_\lambda$ for some $\lambda$
by Lemma \href{https://stacks.math.columbia.edu/tag/0AIA}{lemma factor through thickening (Tag 0AIA)}.
Thus $A \to \Gamma(X_\lambda, \mathcal{O}) \to \Gamma(T, \mathcal{O})$
is surjective. In this way we see that (3) holds.

\medskip\noindent
It is clear that (3) implies (2).

\medskip\noindent
Assume (2) for $A$ and $X \to \Spec(A)$. Write $X = \colim X_\lambda$
as in Definition \ref{definition-affine-formal-algebraic-space}.
Then $A_\lambda = \Gamma(X_\lambda, \mathcal{O})$ is a quotient
of $A$ by assumption (2). Hence $A^\wedge = \lim A_\lambda$
is a complete topological ring, see discussion in
More on Algebra, Section \href{https://stacks.math.columbia.edu/tag/07E7}{section topological ring (Tag 07E7)}.
The maps $A^\wedge \to A_\lambda$ are surjective as $A \to A_\lambda$ is.
We claim that for any $\lambda$ the kernel $I_\lambda \subset A^\wedge$ of
$A^\wedge \to A_\lambda$ is a weak ideal of definition.
Namely, it is open by definition of the limit topology.
If $f \in I_\lambda$, then for any $\mu \in \Lambda$
the image of $f$ in $A_\mu$ is zero in all the residue fields
of the points of $X_\mu$. Hence it is a nilpotent element
of $A_\mu$. Hence some power $f^n \in I_\mu$. Thus $f^n \to 0$
as $n \to 0$. Thus $A^\wedge$ is weakly admissible.
Finally, suppose that $I \subset A^\wedge$ is a weak ideal
of definition. Then $I \subset A^\wedge$ is open and hence there exists
some $\lambda$ such that $I \supset I_\lambda$. Thus we obtain a morphism
$\Spec(A^\wedge/I) \to \Spec(A_\lambda) \to X$.
Then it follows that $X = \colim \Spec(A^\wedge/I)$ where now
the colimit is over all weak ideals of definition.
Thus (1) holds.

\medskip\noindent
Assume (1). In this case it is clear that $X$ is an affine formal
algebraic space. Let $X = \colim X_\lambda$ be any presentation as in
Definition \ref{definition-affine-formal-algebraic-space}.
For each $\lambda$ we can find a weak ideal of definition
$I \subset A$ such that $X_\lambda \to X$ factors through
$\Spec(A/I) \to X$, see Lemma \href{https://stacks.math.columbia.edu/tag/0AIA}{lemma factor through thickening (Tag 0AIA)}.
Then $X_\lambda = \Spec(A/I_\lambda)$ with $I \subset I_\lambda$.
Conversely, for any weak ideal of definition $I \subset A$
the morphism $\Spec(A/I) \to X$ factors through $X_\lambda$
for some $\lambda$, i.e., $I_\lambda \subset I$.
It follows that each $I_\lambda$ is a weak ideal of definition
and that they form a cofinal subset of the set of weak ideals
of definition. Hence $A = \lim A/I = \lim A/I_\lambda$
and we see that (5) is true and moreover that
$A = \lim \Gamma(X_\lambda, \mathcal{O}_{X_\lambda})$.
\end{proof}
\end{document}
